\section{Option-Only Efficient Frontier}

Let $q\in\R^m$ denote option premium exposures per dollar of NAV.  If
$q_i=0.10$, then ten cents of NAV are exposed to option $i$'s return.
Cash is not optimized; it is the accounting numeraire.

\subsection{Unconstrained tangency}

The option-only Sharpe ratio is
\begin{equation}
   \SR(q)=\frac{q^\top\mu_O}{\sqrt{q^\top\Sigma_Oq}},
   \qquad q\ne0.
   \label{eq:sharpe}
\end{equation}

\begin{theorem}[Option-only tangency portfolio]
\label{thm:tangency}
Suppose $\Sigma_O\succ0$ and $\mu_O\ne0$.  Among all nonzero
$q\in\R^m$, every unconstrained maximum-Sharpe portfolio is a nonzero
scalar multiple of
\[
      q^\star = \Sigma_O^{-1}\mu_O .
\]
The maximal squared Sharpe ratio is
$\mu_O^\top\Sigma_O^{-1}\mu_O$.
\end{theorem}

\begin{proof}
Sharpe is homogeneous of degree zero, so impose $q^\top\Sigma_Oq=1$ and
maximize $q^\top\mu_O$.  The Lagrangian is
$L(q,\lambda)=q^\top\mu_O-\lambda(q^\top\Sigma_Oq-1)$.  First-order
conditions give $\mu_O=2\lambda\Sigma_O q$, hence
$q=c\Sigma_O^{-1}\mu_O$.  Choosing $c$ to satisfy
$q^\top\Sigma_Oq=1$ gives the normalized optimizer.  Substituting yields
the maximum value
$\sqrt{\mu_O^\top\Sigma_O^{-1}\mu_O}$.
\end{proof}

\begin{remark}[Equivalence to Markowitz]
Theorem~\ref{thm:tangency} is mathematically the Markowitz tangency
formula.  The difference is not the optimizer.  The difference is the
construction of $\Sigma_O$: stock Markowitz starts from asset returns,
whereas option Markowitz derives option returns from Greek exposure to
underlying, squared-return, and implied-volatility shocks.
\end{remark}

\subsection{NAV and Greek constraints}

Practical option-only portfolios require constraints:
\begin{align}
  \|q\|_1 &\le L, &&\text{gross NAV budget},\\
  \sum_i \max(-q_i,0) &\le L^-, &&\text{short-premium budget},\\
  |q_i| &\le b_i, &&\text{contract concentration},\\
  \sum_{i:u(i)=u}|q_i| &\le L_u, &&\text{underlying concentration},\\
  |g_\Delta^\top q| &\le D, \quad
  |g_\Gamma^\top q|\le G, \quad
  |g_\nu^\top q|\le V, &&\text{Greek budgets}.
\end{align}
Here $g_\Delta$, $g_\Gamma$, and $g_\nu$ are NAV-normalized aggregate
Greek-loading vectors.

With the normalization $\|q\|_1=L$, the constrained problem is
\begin{equation}
\max_q \frac{q^\top\mu_O}{\sqrt{q^\top\Sigma_Oq}}
\quad\text{s.t.}\quad q\in\mathcal{K}_O,
\label{eq:constrained}
\end{equation}
where $\mathcal{K}_O$ is the intersection of the listed linear, absolute
value, and sign constraints.  This is a fractional constrained program;
the implementation supports direct SLSQP solution and uses the closed-form
tangency portfolio as its starting direction.  The empirical runner also reports
closed-form tangency and projected diagnostic portfolios when repeated rolling fits
would otherwise make the paper build impractically slow.

\begin{proposition}[KKT condition]
\label{prop:kkt}
Let $q^\star$ be a differentiable local optimum of
\eqref{eq:constrained} with nonzero return and variance.  Ignoring
inactive constraints, the marginal condition is
\[
 \frac{\mu_O}{\sqrt{v}}
 - \frac{(q^\top\mu_O)\Sigma_Oq}{v^{3/2}}
 =
 A^\top\lambda ,
 \qquad v=q^\top\Sigma_Oq,
\]
where $A^\top\lambda$ is the signed normal cone generated by active NAV,
underlying, and Greek constraints.
\end{proposition}

\begin{proof}
Differentiate \eqref{eq:sharpe}.  The gradient is
$\mu_O/\sqrt{v}-(q^\top\mu_O)\Sigma_Oq/v^{3/2}$.  At a constrained local
optimum this gradient must lie in the normal cone of the feasible set.
For differentiable active constraints, the normal cone is generated by
their gradients, giving the displayed equation.
\end{proof}

\subsection{Estimation}

The empirical estimator follows the standard risk-model route: construct
$B$, estimate $\Omega$ from point-in-time training shocks, estimate
$\Sigma_\varepsilon$ from residual option-bucket returns, shrink the
resulting covariance toward a diagonal target, and repair numerical PSD
drift by clipping negative eigenvalues \citep{ledoitwolf2004}. This is
economically conservative: it does not assume the Greek model explains
all option return risk, and it explicitly leaves a residual covariance
term.
